%!TEX root = ../main.tex
\section{Proof of Theorem \ref{thm:lqr-variance-T-bound}}
\label{sec:app-proof-lqr-variance-T-bound}

\begin{proof}
   Recall the policy gradient estimators
\[
\widehat G_K
=\sum_{t=0}^{T-1}\Sigma^{-1}\varepsilon_t s_t^\top\,R(\tau),
\qquad
\widehat G_{\ell}
=\sum_{t=0}^{T-1}\Big(\Sigma^{-1}(\varepsilon_t\odot \varepsilon_t)-\mathbf 1\Big)\,R(\tau),
\]
and the stacked noise $\bar\varepsilon:=(\varepsilon_0^\top,\dots,\varepsilon_{T-1}^\top)^\top\sim\mathcal N(0,\overline\Sigma)$
with $\overline\Sigma:=I_T\otimes\Sigma$.
   
   Using Cauchy’s inequality on each Frobenius entry, let $\eta_t = \Sigma^{-1}\varepsilon_t$.
   \begin{align*}
   \|\widehat G_K\|_{\mathrm{F}}^2
   &=\Big\|\sum_{t=0}^{T-1}\Sigma^{-1}\varepsilon_t s_t^\top\Big\|_{\mathrm{F}}^2\,R(\tau)^2
   = \sum_{i,j}\Big(\sum_{t=0}^{T-1} \eta_{t,i}s_{t,j}\Big)^2 R(\tau)^2 \\
   &\le \sum_{i,j}
      \Big(\sum_{t=0}^{T-1} \eta_{t,i}^2\Big)
      \Big(\sum_{t=0}^{T-1} s_{t,j}^2\Big) R(\tau)^2
   = \Big(\sum_{t=0}^{T-1} \|\eta_t\|_2^2\Big)
      \Big(\sum_{t=0}^{T-1} \|s_t\|_2^2\Big) R(\tau)^2.
   \end{align*}
   Using stacked $\bar\varepsilon$ and $\bar s$ we obtain
   \begin{equation}
   \label{eq:second-moment-global}
   \VarF(\widehat G_K) \le \mathbb E\big[\|\widehat G_K\|_{\mathrm{F}}^2\big]
   \;\le\;
   \mathbb E\big[\|\overline\Sigma^{-1}\bar\varepsilon\|_2^2\;\|\bar s\|_2^2\;R(\tau)^2\big].
   \end{equation}

   By AM-GM inequality, we have the decomposition
   \begin{equation}
      \VarF(\widehat G_K)
      \;\le\;2\,\underbrace{\mathbb E\!\big[\|\overline\Sigma^{-1}\bar\varepsilon\|_2^2\,
                                 \|\mathcal F_S s_0\|_2^2\,R(\tau)^{2}\big]}_{(\star)}
      \;+\;2\,\underbrace{\mathbb E\!\big[\|\overline\Sigma^{-1}\bar\varepsilon\|_2^2\,
                                 \|\mathcal F_E\bar\varepsilon\|_2^2\,R(\tau)^{2}\big]}_{(\dagger)} .
    \end{equation}

   For the $\ell$ part, define the per-step score term
\[
g_t \;:=\;\Sigma^{-1}(\varepsilon_t\odot \varepsilon_t)-\mathbf 1\in\R^m,
\qquad
\widehat G_\ell \;=\;\sum_{t=0}^{T-1} g_t\,R(\tau).
\]
By Cauchy--Schwarz (or $\|\sum_{t=0}^{T-1}v_t\|_2^2\le T\sum_{t=0}^{T-1}\|v_t\|_2^2$),
\begin{equation}\label{eq:ell-second-moment-global}
\mathbb E\big[\|\widehat G_\ell\|_2^2\big]
\;\le\;
T\,\mathbb E\big[\sum_{t=0}^{T-1}\|g_t\|_2^2\,R(\tau)^2\big].
\end{equation}

We next bound $\sum_{t=0}^{T-1}\|g_t\|_2^2$ by a polynomial in $\bar\varepsilon$.
Since $\Sigma$ is diagonal, write $\varepsilon_t=\Sigma^{1/2}\xi_t$ with $\xi_t\sim\mathcal N(0,I_m)$ i.i.d.,
so that
\[
g_t=\Sigma^{-1}(\varepsilon_t\odot \varepsilon_t)-\mathbf 1
= (\xi_t\odot\xi_t)-\mathbf 1,
\qquad
\|g_t\|_2^2=\sum_{i=1}^m(\xi_{t,i}^2-1)^2.
\]
Using $(a-1)^2\le 2a^2+2$ with $a=\xi_{t,i}^2\ge 0$ gives
\[
(\xi_{t,i}^2-1)^2 \le 2\xi_{t,i}^4+2
\quad\Rightarrow\quad
\|g_t\|_2^2 \le 2\sum_{i=1}^m \xi_{t,i}^4 + 2m
\le 2\|\xi_t\|_2^4 + 2m.
\]
Summing over $t$ yields
\[
\sum_{t=0}^{T-1}\|g_t\|_2^2
\le 2\sum_{t=0}^{T-1}\|\xi_t\|_2^4 + 2mT.
\]
Using $\|\xi_t\|_2^2\ge 0$, we have 
$
\sum_{t=0}^{T-1}\|\xi_t\|_2^4
\le \Big(\sum_{t=0}^{T-1}\|\xi_t\|_2^2\Big)^2
=\|\bar\xi\|_2^4.
$

Finally, since $\bar\varepsilon=(I_T\otimes\Sigma^{1/2})\bar\xi$ and $\overline\Sigma=I_T\otimes\Sigma$,
we have $\overline\Sigma^{-1/2}\bar\varepsilon=\bar\xi$, so $\|\bar\xi\|_2^4=\|\overline\Sigma^{-1/2}\bar\varepsilon\|_2^4$.
Therefore,
\[
\sum_{t=0}^{T-1}\|g_t\|_2^2
\le 2\|\overline\Sigma^{-1/2}\bar\varepsilon\|_2^4 + 2mT.
\]
Plugging into \eqref{eq:ell-second-moment-global} gives
\[
\mathbb E\big[\|\widehat G_\ell\|_2^2\big]
\le
T\,\mathbb E\!\Big[\big(2\|\overline\Sigma^{-1/2}\bar\varepsilon\|_2^4 + 2mT\big)\,R(\tau)^2\Big].
\]


Assume the standard multi-step quadratic decomposition of the return
\[
R(\tau)=-(x+2y+z),
\qquad
x:=s_0^\top \overline M_{ss}s_0,\quad
y:=s_0^\top \overline M_{se}\bar\varepsilon,\quad
z:=\bar\varepsilon^\top \overline M_{ee}\bar\varepsilon,
\]
\paragraph{Proof of $K$-part expansion.}
Define
\[
S:=\mathcal F_S^\top\mathcal F_S,\qquad
E:=\mathcal F_E^\top\mathcal F_E,\qquad
\]


\paragraph{Step 1: expand $(\star)$.}
Since $\|\mathcal F_S s_0\|_2^2=s_0^\top S s_0$ and $R(\tau)=-(x+2y+z)$,
\[
(\star)
=
\E\!\left[(\bar\varepsilon^\top  \overline\Sigma^{-2}\bar\varepsilon)\,(s_0^\top S s_0)\,(x+2y+z)^2\right].
\]
Expanding $(x+2y+z)^2=x^2+z^2+2xz+4y^2+4xy+4yz$ and using independence
$s_0\perp\!\!\!\perp\bar\varepsilon$ and odd-moment cancellation in $\bar\varepsilon$
(which removes the $xy$ and $yz$ terms), we get
\[
(\star)=E^{(\star)}_{1,K}+E^{(\star)}_{2,K}+E^{(\star)}_{3,K}+E^{(\star)}_{4,K},
\]
with
\begin{align*}
E^{(\star)}_{1,K}
&:= \E[\bar\varepsilon^\top  \overline\Sigma^{-2}\bar\varepsilon]\;\E\!\left[(s_0^\top S s_0)\,x^2\right],\\
E^{(\star)}_{2,K}
&:= \E[s_0^\top S s_0]\;\E\!\left[(\bar\varepsilon^\top  \overline\Sigma^{-2}\bar\varepsilon)\,z^2\right],\\
E^{(\star)}_{3,K}
&:= 2\,\E\!\left[(s_0^\top S s_0)\,x\right]\;\E\!\left[(\bar\varepsilon^\top  \overline\Sigma^{-2}\bar\varepsilon)\,z\right],\\
E^{(\star)}_{4,K}
&:= 4\,\E\!\left[(s_0^\top S s_0)\;
\E\!\left[(\bar\varepsilon^\top  \overline\Sigma^{-2}\bar\varepsilon)\,y^2\,\middle|\,s_0\right]\right].
\end{align*}

Using Lemma~\ref{lem:IS_shorthand},
\[
\E[\bar\varepsilon^\top  \overline\Sigma^{-2}\bar\varepsilon]=\mathrm{IS}_{\overline\Sigma}( \overline\Sigma^{-2}),
\qquad
\E[s_0^\top S s_0]=\mathrm{IS}_{\Sigma_0}(S),
\]
\[
\E\!\left[(s_0^\top S s_0)x\right]=\mathrm{IS}_{\Sigma_0}(S,\overline M_{ss}),
\qquad
\E\!\left[(s_0^\top S s_0)x^2\right]=\mathrm{IS}_{\Sigma_0}(S,\overline M_{ss},\overline M_{ss}),
\]
and since $z=\bar\varepsilon^\top \overline M_{ee}\bar\varepsilon$,
\[
\E\!\left[(\bar\varepsilon^\top  \overline\Sigma^{-2}\bar\varepsilon)z\right]
=\mathrm{IS}_{\overline\Sigma}( \overline\Sigma^{-2},\overline M_{ee}),
\qquad
\E\!\left[(\bar\varepsilon^\top  \overline\Sigma^{-2}\bar\varepsilon)z^2\right]
=\mathrm{IS}_{\overline\Sigma}( \overline\Sigma^{-2},\overline M_{ee},\overline M_{ee}).
\]
For $E^{(\star)}_{4,K}$, note $y=s_0^\top \overline M_{se}\bar\varepsilon$ implies
\[
y^2=\bar\varepsilon^\top X(s_0)\bar\varepsilon,
\qquad
X(s_0):=\overline M_{se}^\top s_0 s_0^\top \overline M_{se}.
\]
Conditioning on $s_0$ and applying Lemma~\ref{lem:IS_shorthand} with $k=2$,
\[
\E_{\bar\varepsilon}\!\left[(\bar\varepsilon^\top  \overline\Sigma^{-2}\bar\varepsilon)y^2\mid s_0\right]
=\mathrm{IS}_{\overline\Sigma}\!\big( \overline\Sigma^{-2},X(s_0)\big).
\]
Define
\[
U:=\overline M_{se}\overline M_{se}^\top,
\qquad
V:=\overline M_{se}\overline\Sigma \overline M_{se}^\top.
\]
Then $\tr(X(s_0))=s_0^\top U s_0$ and $\tr(\overline\Sigma X(s_0))=s_0^\top V s_0$, and since
$\overline\Sigma  \overline\Sigma^{-2}\overline\Sigma=I$, we obtain
\[
\mathrm{IS}_{\overline\Sigma}\!\big( \overline\Sigma^{-2},X(s_0)\big)
=
\mathrm{IS}_{\overline\Sigma}( \overline\Sigma^{-2})\,(s_0^\top V s_0)+2\,(s_0^\top U s_0).
\]
Therefore,
\[
E^{(\star)}_{4,K}
=
4\Big(
\mathrm{IS}_{\overline\Sigma}( \overline\Sigma^{-2})\;\mathrm{IS}_{\Sigma_0}(S,V)
+
2\,\mathrm{IS}_{\Sigma_0}(S,U)
\Big).
\]
Altogether,
\begin{equation}\label{eq:star-final}
\begin{aligned}
(\star)
&=
\mathrm{IS}_{\overline\Sigma}( \overline\Sigma^{-2})\;\mathrm{IS}_{\Sigma_0}(S,\overline M_{ss},\overline M_{ss})
+\mathrm{IS}_{\Sigma_0}(S)\;\mathrm{IS}_{\overline\Sigma}( \overline\Sigma^{-2},\overline M_{ee},\overline M_{ee})\\
&\quad
+2\,\mathrm{IS}_{\Sigma_0}(S,\overline M_{ss})\;\mathrm{IS}_{\overline\Sigma}( \overline\Sigma^{-2},\overline M_{ee})
+4\Big(
\mathrm{IS}_{\overline\Sigma}( \overline\Sigma^{-2})\;\mathrm{IS}_{\Sigma_0}(S,V)
+2\,\mathrm{IS}_{\Sigma_0}(S,U)
\Big).
\end{aligned}
\end{equation}

\paragraph{Step 2: expand $(\dagger)$.}
Since $\|\mathcal F_E\bar\varepsilon\|_2^2=\bar\varepsilon^\top E\bar\varepsilon$,
\[
(\dagger)
=
\E\!\left[(\bar\varepsilon^\top  \overline\Sigma^{-2}\bar\varepsilon)\,(\bar\varepsilon^\top E\bar\varepsilon)\,(x+2y+z)^2\right].
\]
As before, odd-moment cancellation removes $xy$ and $yz$, yielding
\[
(\dagger)=E^{(\dagger)}_{1,K}+E^{(\dagger)}_{2,K}+E^{(\dagger)}_{3,K}+E^{(\dagger)}_{4,K},
\]
with
\begin{align*}
E^{(\dagger)}_{1,K}
&:= \E\!\left[(\bar\varepsilon^\top  \overline\Sigma^{-2}\bar\varepsilon)(\bar\varepsilon^\top E\bar\varepsilon)\right]\;\E[x^2],\\
E^{(\dagger)}_{2,K}
&:= \E\!\left[(\bar\varepsilon^\top  \overline\Sigma^{-2}\bar\varepsilon)(\bar\varepsilon^\top E\bar\varepsilon)\,z^2\right],\\
E^{(\dagger)}_{3,K}
&:= 2\,\E[x]\;\E\!\left[(\bar\varepsilon^\top  \overline\Sigma^{-2}\bar\varepsilon)(\bar\varepsilon^\top E\bar\varepsilon)\,z\right],\\
E^{(\dagger)}_{4,K}
&:= 4\,\E\!\left[\E\!\left[(\bar\varepsilon^\top  \overline\Sigma^{-2}\bar\varepsilon)(\bar\varepsilon^\top E\bar\varepsilon)\,y^2\,\middle|\,s_0\right]\right].
\end{align*}
Rewrite each term:
\[
\E[x]=\mathrm{IS}_{\Sigma_0}(\overline M_{ss}),
\qquad
\E[x^2]=\mathrm{IS}_{\Sigma_0}(\overline M_{ss},\overline M_{ss}),
\]
\[
\E\!\left[(\bar\varepsilon^\top  \overline\Sigma^{-2}\bar\varepsilon)(\bar\varepsilon^\top E\bar\varepsilon)\right]
=\mathrm{IS}_{\overline\Sigma}( \overline\Sigma^{-2},E),
\]
\[
\E\!\left[(\bar\varepsilon^\top  \overline\Sigma^{-2}\bar\varepsilon)(\bar\varepsilon^\top E\bar\varepsilon)\,z\right]
=\mathrm{IS}_{\overline\Sigma}( \overline\Sigma^{-2},E,\overline M_{ee}),
\qquad
\E\!\left[(\bar\varepsilon^\top  \overline\Sigma^{-2}\bar\varepsilon)(\bar\varepsilon^\top E\bar\varepsilon)\,z^2\right]
=\mathrm{IS}_{\overline\Sigma}( \overline\Sigma^{-2},E,\overline M_{ee},\overline M_{ee}).
\]
For the mixed term, with $y^2=\bar\varepsilon^\top X(s_0)\bar\varepsilon$,
\[
\E_{\bar\varepsilon}\!\left[(\bar\varepsilon^\top  \overline\Sigma^{-2}\bar\varepsilon)(\bar\varepsilon^\top E\bar\varepsilon)\,y^2\mid s_0\right]
=\mathrm{IS}_{\overline\Sigma}\!\big( \overline\Sigma^{-2},E,X(s_0)\big),
\]
For the mixed term, recall $X(s_0)=\overline M_{se}^\top s_0 s_0^\top \overline M_{se}$ and define
\[
U:=\overline M_{se}\overline M_{se}^\top,\qquad
V:=\overline M_{se}\overline\Sigma\,\overline M_{se}^\top,\qquad
W:=\overline M_{se}\,\overline\Sigma\,E\,\overline\Sigma\,\overline M_{se}^\top,\qquad
Z:=\symfun(\overline M_{se}\,E\,\overline\Sigma\,\overline M_{se}^\top).
\]
First, by Lemma~\ref{lem:IS_shorthand} ($k=3$) with $\Omega=\overline\Sigma$, $A=\overline\Sigma^{-2}$, $B=E$, $C=X(s_0)$,
\begin{align*}
\mathrm{IS}_{\overline\Sigma}(\overline\Sigma^{-2},E,X(s_0))
&=\tr(\overline\Sigma^{-1})\tr(\overline\Sigma E)\tr(\overline\Sigma X(s_0))
+2\tr(E)\tr(\overline\Sigma X(s_0))
+2\tr(\overline\Sigma E)\tr(X(s_0))\\
&\quad+2\tr(\overline\Sigma^{-1})\tr(\overline\Sigma E\overline\Sigma X(s_0))
+8\tr(E\overline\Sigma X(s_0)),
\end{align*}
and using $\overline\Sigma\,\overline\Sigma^{-2}\,\overline\Sigma=I$ together with
$\tr(\overline\Sigma X(s_0))=s_0^\top V s_0$, $\tr(X(s_0))=s_0^\top U s_0$,
$\tr(\overline\Sigma E\overline\Sigma X(s_0))=s_0^\top W s_0$, and $\tr(E\overline\Sigma X(s_0))=s_0^\top Z s_0$,
we obtain
\[
\mathrm{IS}_{\overline\Sigma}(\overline\Sigma^{-2},E,X(s_0))
=
\mathrm{IS}_{\overline\Sigma}(\overline\Sigma^{-2},E)\,(s_0^\top V s_0)
+2\,\mathrm{IS}_{\overline\Sigma}(E)\,(s_0^\top U s_0)
+2\,\mathrm{IS}_{\overline\Sigma}(\overline\Sigma^{-2})\,(s_0^\top W s_0)
+8\,(s_0^\top Z s_0).
\]

Taking expectation over $s_0\sim\mathcal N(0,\Sigma_0)$ gives
\[
\E_{s_0}\!\left[\mathrm{IS}_{\overline\Sigma}(\overline\Sigma^{-2},E,X(s_0))\right]
=
\mathrm{IS}_{\overline\Sigma}(\overline\Sigma^{-2},E)\,\mathrm{IS}_{\Sigma_0}(V)
+2\,\mathrm{IS}_{\overline\Sigma}(E)\,\mathrm{IS}_{\Sigma_0}(U)
+2\,\mathrm{IS}_{\overline\Sigma}(\overline\Sigma^{-2})\,\mathrm{IS}_{\Sigma_0}(W)
+8\,\mathrm{IS}_{\Sigma_0}(Z).
\]
Consequently,
\[
E^{(\dagger)}_{4,K}
=
4\,\mathrm{IS}_{\overline\Sigma}(\overline\Sigma^{-2},E)\,\mathrm{IS}_{\Sigma_0}(V)
+8\,\mathrm{IS}_{\overline\Sigma}(E)\,\mathrm{IS}_{\Sigma_0}(U)
+8\,\mathrm{IS}_{\overline\Sigma}(\overline\Sigma^{-2})\,\mathrm{IS}_{\Sigma_0}(W)
+32\,\mathrm{IS}_{\Sigma_0}(Z).
\]

hence
\begin{equation}\label{eq:dagger-final}
\begin{aligned}
(\dagger)
&=
\mathrm{IS}_{\overline\Sigma}( \overline\Sigma^{-2},E)\;\mathrm{IS}_{\Sigma_0}(\overline M_{ss},\overline M_{ss})
+\mathrm{IS}_{\overline\Sigma}( \overline\Sigma^{-2},E,\overline M_{ee},\overline M_{ee})\\
&\quad
+2\,\mathrm{IS}_{\Sigma_0}(\overline M_{ss})\;\mathrm{IS}_{\overline\Sigma}( \overline\Sigma^{-2},E,\overline M_{ee})
+4\,\E_{s_0}\!\left[\mathrm{IS}_{\overline\Sigma}\!\big( \overline\Sigma^{-2},E,X(s_0)\big)\right].
\end{aligned}
\end{equation}

Combining \eqref{eq:star-final}, and \eqref{eq:dagger-final} yields the desired
multi-step expansion of $\E[\|\widehat G_K\|_{\mathrm{F}}^2]$.

% ============================================================
\paragraph{Proof of $\ell$-part expansion.}
Then
\[
\|\overline\Sigma^{-1/2}\bar\varepsilon\|_2^4
=
(\bar\varepsilon^\top \overline\Sigma^{-1} \bar\varepsilon)^2,
\qquad
R(\tau)^2=(x+2y+z)^2.
\]
As above, odd-moment cancellation removes $xy$ and $yz$, so
\[
R(\tau)^2=x^2+z^2+2xz+4y^2.
\]
Therefore
\[
\E\!\left[(2\|\overline\Sigma^{-1/2}\bar\varepsilon\|_2^4+2mT)\,R(\tau)^2\right]
=
2\sum_{i=1}^4 A_i + 2mT\sum_{i=1}^4 B_i,
\]
where
\begin{align*}
A_1 &:= \E\!\left[(\bar\varepsilon^\top \overline\Sigma^{-1}\bar\varepsilon)^2\,x^2\right]
     = \E[x^2]\;\E\!\left[(\bar\varepsilon^\top \overline\Sigma^{-1}\bar\varepsilon)^2\right],\\
A_2 &:= \E\!\left[(\bar\varepsilon^\top \overline\Sigma^{-1}\bar\varepsilon)^2\,z^2\right],\\
A_3 &:= 2\,\E\!\left[(\bar\varepsilon^\top \overline\Sigma^{-1}\bar\varepsilon)^2\,x z\right]
     = 2\,\E[x]\;\E\!\left[(\bar\varepsilon^\top \overline\Sigma^{-1}\bar\varepsilon)^2\,z\right],\\
A_4 &:= 4\,\E\!\left[(\bar\varepsilon^\top \overline\Sigma^{-1}\bar\varepsilon)^2\,y^2\right]
     = 4\,\E_{s_0}\!\left[\mathrm{IS}_{\overline\Sigma}(\overline\Sigma^{-1},\overline\Sigma^{-1},X(s_0))\right],
\end{align*}
and
\begin{align*}
B_1 &:= \E[x^2],\qquad
B_2 := \E[z^2],\qquad
B_3 := 2\,\E[xz]=2\,\E[x]\E[z],\qquad
B_4 := 4\,\E[y^2].
\end{align*}

Each term can be written using $\mathrm{IS}$ shorthand:
\[
\E[x]=\mathrm{IS}_{\Sigma_0}(\overline M_{ss}),
\qquad
\E[x^2]=\mathrm{IS}_{\Sigma_0}(\overline M_{ss},\overline M_{ss}),
\]
\[
\E[z]=\mathrm{IS}_{\overline\Sigma}(\overline M_{ee}),
\qquad
\E[z^2]=\mathrm{IS}_{\overline\Sigma}(\overline M_{ee},\overline M_{ee}),
\]
\[
\E[y^2]=\mathrm{IS}_{\Sigma_0}(V),
\quad\text{where }V:=\overline M_{se}\overline\Sigma \overline M_{se}^\top.
\]
Moreover,
\[
\E\!\left[(\bar\varepsilon^\top \overline\Sigma^{-1}\bar\varepsilon)^2\right]=\mathrm{IS}_{\overline\Sigma}(\overline\Sigma^{-1},\overline\Sigma^{-1}),
\]
\[
\E\!\left[(\bar\varepsilon^\top \overline\Sigma^{-1}\bar\varepsilon)^2\,z\right]
=\mathrm{IS}_{\overline\Sigma}(\overline\Sigma^{-1},\overline\Sigma^{-1},\overline M_{ee}),
\qquad
\E\!\left[(\bar\varepsilon^\top \overline\Sigma^{-1}\bar\varepsilon)^2\,z^2\right]
=\mathrm{IS}_{\overline\Sigma}(\overline\Sigma^{-1},\overline\Sigma^{-1},\overline M_{ee},\overline M_{ee}).
\]

\noindent\textbf{Term $A_4$ ($y^2$-part).}
Applying Lemma~\ref{lem:IS_shorthand} ($k=3$) with $\Omega=\overline\Sigma$, $A=B=\overline\Sigma^{-1}$, $C=X(s_0)$ and using
$\overline\Sigma\,\overline\Sigma^{-1}=I$, we obtain
\[
\mathrm{IS}_{\overline\Sigma}(\overline\Sigma^{-1},\overline\Sigma^{-1},X(s_0))
=\big(d^2+6d+8\big)\,\tr\!\big(\overline\Sigma\,X(s_0)\big)
=(d+2)(d+4)\,\tr\!\big(\overline\Sigma\,X(s_0)\big),
\qquad d:=mT.
\]
Moreover,
\[
\tr(\overline\Sigma\,X(s_0))
=\tr\!\big(s_0 s_0^\top \overline M_{se}\overline\Sigma\,\overline M_{se}^\top\big)
=s_0^\top V s_0,
\qquad
V:=\overline M_{se}\overline\Sigma\,\overline M_{se}^\top,
\]
hence $\E_{s_0}[\tr(\overline\Sigma\,X(s_0))]=\tr(\Sigma_0 V)=\mathrm{IS}_{\Sigma_0}(V)$.
Therefore,
\[
A_4
=4(d+2)(d+4)\,\mathrm{IS}_{\Sigma_0}(V),
\qquad d=mT.
\]

Hence,
\begin{equation}\label{eq:Ell-final}
\begin{aligned}
\E\!\left[\|\widehat G_\ell\|_2^2\right]
&\le
T\Bigg(
2\Big[
\mathrm{IS}_{\Sigma_0}(\overline M_{ss},\overline M_{ss})\;\mathrm{IS}_{\overline\Sigma}(\overline\Sigma^{-1},\overline\Sigma^{-1})
+\mathrm{IS}_{\overline\Sigma}(\overline\Sigma^{-1},\overline\Sigma^{-1},\overline M_{ee},\overline M_{ee})\\
&\qquad\qquad
+2\,\mathrm{IS}_{\Sigma_0}(\overline M_{ss})\;\mathrm{IS}_{\overline\Sigma}(\overline\Sigma^{-1},\overline\Sigma^{-1},\overline M_{ee})
+4(d+2)(d+4)\,\mathrm{IS}_{\Sigma_0}(V)
\Big]\\
&\qquad
+2mT\Big[
\mathrm{IS}_{\Sigma_0}(\overline M_{ss},\overline M_{ss})
+\mathrm{IS}_{\overline\Sigma}(\overline M_{ee},\overline M_{ee})
+2\,\mathrm{IS}_{\Sigma_0}(\overline M_{ss})\,\mathrm{IS}_{\overline\Sigma}(\overline M_{ee})\\
&\qquad\qquad
+4\,\mathrm{IS}_{\Sigma_0}(V)
\Big]
\Bigg),
\end{aligned}
\end{equation}
where $X(s_0):=\overline M_{se}^\top s_0 s_0^\top \overline M_{se}$ and $V:=\overline M_{se}\overline\Sigma \overline M_{se}^\top$.

\end{proof}


\section{Proof of Theorem \ref{thm:lqr-lifted-state-maps-bound}}
\label{sec:app-proof-lqr-lifted-state-maps-bound}

\begin{proof}
Recall that $\|\mathcal F_S\|_2^2=\lambda_{\max}(\mathcal F_S^\top \mathcal F_S)$. Since
\[
\mathcal F_S=\begin{bmatrix}I & F & \cdots & F^{T-1}\end{bmatrix}^\top,
\]
we have
\[
\mathcal F_S^\top \mathcal F_S=\sum_{t=0}^{T-1}(F^t)^\top F^t.
\]

\paragraph{Lower bound.}
For any fixed $t\in\{0,\dots,T-1\}$,
\[
\mathcal F_S^\top \mathcal F_S
=\sum_{k=0}^{T-1}(F^k)^\top F^k
\succeq (F^t)^\top F^t.
\]
Taking $\lambda_{\max}(\cdot)$ preserves the Loewner order, hence
\[
\|\mathcal F_S\|_2^2
=\lambda_{\max}(\mathcal F_S^\top \mathcal F_S)
\ge \lambda_{\max}\big((F^t)^\top F^t\big)
=\|F^t\|_2^2.
\]
Maximizing over $t$ yields
$\max_{0\le t\le T-1}\|F^t\|_2^2 \le \|\mathcal F_S\|_2^2$.

\paragraph{Upper bound.}
For any unit vector $v$,
\[
v^\top(\mathcal F_S^\top \mathcal F_S)v
=\sum_{t=0}^{T-1}\|F^t v\|_2^2
\le \sum_{t=0}^{T-1}\|F^t\|_2^2\,\|v\|_2^2
=\sum_{t=0}^{T-1}\|F^t\|_2^2.
\]
Taking the supremum over $\|v\|_2=1$ gives
\[
\|\mathcal F_S\|_2^2
=\sup_{\|v\|_2=1} v^\top(\mathcal F_S^\top \mathcal F_S)v
\le \sum_{t=0}^{T-1}\|F^t\|_2^2.
\]
\end{proof}


